\chapter*{Abstract}
\addcontentsline{toc}{chapter}{Abstract}
\markboth{ABSTRACT}{ABSTRACT}

Ion energy flux across the plasma sheath governs the surface heating, erosion, and
efficiency loss that ultimately limit the performance and lifetime of every
plasma--surface technology, from electric-propulsion thrusters and semiconductor
plasma-etch reactors to the plasma-facing components of fusion devices. Although
first-principles fluid and kinetic models can \emph{predict} this flux from the
sheath potential and the local plasma parameters, the predictions remain weakly
validated under realistic, non-equilibrium, and time-varying conditions. The
measurement tools available today fall short in complementary ways: intrusive
electrostatic probes perturb the very sheath they sample and cannot survive
high-flux transient environments, while individual non-invasive optical
techniques capture only a partial slice of the plasma state and are rarely
time-resolved. As a consequence, no widely accepted, experimentally validated
framework exists for quantifying ion energy transfer at plasma--surface
boundaries.

This monograph develops, from first principles, a non-intrusive diagnostic
framework that reconstructs ion energy flux by posing it as a physics-constrained
inverse problem. Four mature, mutually complementary optical diagnostics --
optical emission spectroscopy, laser-induced fluorescence, Thomson scattering,
and interferometry -- are fused to over-determine the plasma state without
inserting any object into the plasma. The resulting redundant observables are
regularized with sheath theory and kinetic simulation, transforming partial,
indirect, and noisy measurements into a stable, spatially and temporally resolved
estimate of the energy deposited at the wall.

The document is organised as a complete reference work. Part~I establishes the
physics of plasmas, sheaths, ion energy flux, and plasma--surface interaction.
Part~II surveys the measurement and modelling landscape in depth -- intrusive and
optical diagnostics, fluid and kinetic models, inverse-problem theory, and the
proposed physics-constrained inversion. Part~III consolidates a structured
review of the supporting literature and sets out open problems and conclusions. Throughout, the
treatment is self-contained and pedagogically ordered, so that a reader with no
prior plasma background can follow the argument from the definition of the fourth
state of matter to the design of a deployable, validated ion-energy-flux
diagnostic.

\vspace{0.5cm}
\noindent\textbf{Keywords:} plasma sheath; ion energy flux; non-intrusive
diagnostics; optical emission spectroscopy; laser-induced fluorescence; Thomson
scattering; interferometry; inverse problems; physics-constrained inversion;
electric propulsion; plasma etching; fusion plasma-facing components.
\cleardoublepage
